\begin{tikzpicture}[font=\small,line cap=round,line join=round]\begin{scope}[x=1.5cm,y=0.4444444444444444cm]\draw[otRule!65,line width=.2pt](0,0)--(0,9);\draw[otRule!65,line width=.2pt](1,0)--(1,9);\node[below,font=\tiny,fill=white,inner sep=.5pt]at(1,0){1};\draw[otRule!65,line width=.2pt](2,0)--(2,9);\node[below,font=\tiny,fill=white,inner sep=.5pt]at(2,0){2};\draw[otRule!65,line width=.2pt](3,0)--(3,9);\node[below,font=\tiny,fill=white,inner sep=.5pt]at(3,0){3};\draw[otRule!65,line width=.2pt](4,0)--(4,9);\node[below,font=\tiny,fill=white,inner sep=.5pt]at(4,0){4};\draw[otRule!65,line width=.2pt](5,0)--(5,9);\node[below,font=\tiny,fill=white,inner sep=.5pt]at(5,0){5};\draw[otRule!65,line width=.2pt](6,0)--(6,9);\node[below,font=\tiny,fill=white,inner sep=.5pt]at(6,0){6};\draw[otRule!65,line width=.2pt](0,0.0)--(6,0.0);\draw[otRule!65,line width=.2pt](0,1.5)--(6,1.5);\node[left,font=\tiny,fill=white,inner sep=.5pt]at(0,1.5){1.5};\draw[otRule!65,line width=.2pt](0,3.0)--(6,3.0);\node[left,font=\tiny,fill=white,inner sep=.5pt]at(0,3.0){3.0};\draw[otRule!65,line width=.2pt](0,4.5)--(6,4.5);\node[left,font=\tiny,fill=white,inner sep=.5pt]at(0,4.5){4.5};\draw[otRule!65,line width=.2pt](0,6.0)--(6,6.0);\node[left,font=\tiny,fill=white,inner sep=.5pt]at(0,6.0){6.0};\draw[otRule!65,line width=.2pt](0,7.5)--(6,7.5);\node[left,font=\tiny,fill=white,inner sep=.5pt]at(0,7.5){7.5};\draw[otRule!65,line width=.2pt](0,9.0)--(6,9.0);\node[left,font=\tiny,fill=white,inner sep=.5pt]at(0,9.0){9.0};\draw[otMuted,->](0,0)--(6.1,0)node[right]{m};\draw[otMuted,->](0,0)--(0,9.45)node[above]{euros};\node[below left,font=\tiny]at(0,0){0};\begin{scope}\clip(0,0)rectangle(6,9);\geoSolution{\draw[otNavy,line width=1pt]plot[domain=0:6,samples=90](\x,{1.5*\x});}\geoSolution{}\end{scope}\end{scope}\end{tikzpicture}\par\vspace{1mm}{\fontsize{8}{10}\selectfont\color{otMuted} Les axes sont gradués ; leurs échelles peuvent différer.}